\section{策略安全与正则化}

\subsection{防止策略崩溃的 Regularization}

无论是 PPO 还是 SAC，当策略更新过猛都会导致性能崩盘。实务中常用的 regularization 包括：

\begin{itemize}
    \item \textbf{KL penalty}：在 objective 中加入 $\lambda D_{\mathrm{KL}}(\pi_{\theta'}\|\pi_\theta)$，直接惩罚过大偏移。
    \item \textbf{Entropy decay}：在 SAC 中调节 $\alpha$，让策略在训练后期逐步减少探索以收敛。
    \item \textbf{Gradient clipping}：限制梯度范数避免步长过大。
\end{itemize}

\subsection{监控策略行为}

推荐的监控指标：

\begin{itemize}
    \item \textbf{Average KL divergence} between successive policies
    \item \textbf{Episode variance} of returns—high variance signals instability
    \item \textbf{Q-value rollout test}—simulate fixed rollout to check escalation
\end{itemize}

\begin{knowledgebox}{自动告警设计}
设定门限（例如 KL > 0.05 时报错，episode std dev 增加 30% 时停止训练）可以让 pipeline 自动回退到 prior checkpoint。这种机制在 Walls-Free RL（Sim2Real）项目中广泛采用。
\end{knowledgebox}

\subsection{本章小结}

安全调度需要 regularization 与监控指标结合。KL penalty、entropy decay、gradient clipping 是常见防过度偏移手段，而 KL、variance 和 Q-value rollout 监控提供实时风险预警。

